% !TEX root = ../main.tex
\chapter{Experiment Results}

\section{Telemetry Mapping and Analysis Assumptions}

The results in this chapter come from
\path{telemetry_logs/20260411_142117/telemetry.csv}. The analysis script uses the
logged columns \verb|t_ms|, \verb|setpoint_cm|, \verb|d_filt_cm|, \verb|x_cm|,
\verb|x_dot_cm_s|, \verb|theta_rel_deg|, \verb|theta_dot_deg_s|,
\verb|theta_cmd_rel_deg|, \verb|mode|, \verb|valid|, \verb|driver_enabled|,
\verb|recovery_active|, \verb|step_pos|, and \verb|step_delta|.

Only rows with \verb|mode == 2| and \verb|driver_enabled == 1| are treated as
closed-loop experiment data. Position-derived metrics also honor the logged
\verb|valid| flag because the far-step segment contains intermittent invalid Sharp
samples.

For figure readability, the position and error plots use a display-only five-sample
centered rolling mean of the logged position signal, and short interior invalid gaps
are bridged only in the drawn trace. The summary tables and reported metrics still use
the raw valid samples from the CSV.

\ReportTableFromFile
  {Signal mapping assumptions used by the analysis pipeline.}
  {tables/signal_mapping.tex}
  {tab:signal_mapping}

\section{Run Summary}

The full log contains \TelemetryRowCount{} rows spanning \TelemetryDurationSeconds{}
seconds, with \TelemetryClosedLoopRowCount{} rows in closed-loop cascade mode. The run
contains \TelemetryInvalidRowCount{} invalid sensor rows and \TelemetryRecoveryRowCount{}
rows with recovery active. The maximum observed absolute tracking error was
\TelemetryMaxAbsErrorCm{} cm, while the maximum measured beam-angle magnitude and
command magnitude were \TelemetryMaxBeamAngleDeg{} deg and \TelemetryMaxCommandDeg{}
deg respectively.

\ReportTableFromFile
  {Run-level summary generated from the telemetry CSV.}
  {tables/telemetry_run_summary.tex}
  {tab:telemetry_run_summary}

\section{Overview Plot}

\ReportFigure[width=\linewidth]
  {figures/telemetry_overview.pdf}
  {Phase-by-phase overview of the staircase experiment. Each panel uses a local time
  axis and its own vertical scale so the hold portions remain readable. The bold red
  curve is a display-smoothed version of the valid measured signal, while the faint raw
  trace shows the underlying sensor fluctuation.}
  {fig:telemetry_overview}

The staircase sequence consists of an initial center hold, a far-side step, a return to
center, and a near-side step. The local time axes make the settled portions legible
without altering the underlying data. Validity losses are concentrated in the far-side
segment, which is why the generated summary tables separate valid sample counts from
total row counts.

\section{Ball-Position Tracking}

\ReportFigure[width=\linewidth]
  {figures/telemetry_position_tracking.pdf}
  {Reference tracking by commanded transition. The left column shows the first
  \SI{8}{s} after each setpoint change, while the right column zooms into the late hold
  portion of the same segment.}
  {fig:telemetry_position_tracking}

\ReportFigure[width=\linewidth]
  {figures/telemetry_tracking_error.pdf}
  {Late-hold tracking error
  $x_{\text{cm}} = d_{\text{filt,cm}} - d_{\text{setpoint,cm}}$ for each commanded
  transition. Dashed lines show the default settling band used in the summary metrics,
  while the bold curve uses the same display-only smoothing as the position plots.}
  {fig:telemetry_tracking_error}

The split transition-and-hold layout makes the experiment easier to read than a single
long time axis. In particular, the far-side and near-side moves show usable late-hold
tracking once the initial transport transient is separated from the rest of the phase.
The center-return segment still carries the largest residual offset in the available
window, so the figure presentation is clearer but not artificially optimistic.

\section{Beam Motion and Control Effort}

\ReportFigure[width=\linewidth]
  {figures/telemetry_beam_angle.pdf}
  {Measured beam angle versus time during the staircase experiment.}
  {fig:telemetry_beam_angle}

\ReportFigure[width=\linewidth]
  {figures/telemetry_control_input.pdf}
  {Final relative beam-angle command sent to the inner loop.}
  {fig:telemetry_control_input}

\ReportFigure[width=\linewidth]
  {figures/telemetry_velocity_rate.pdf}
  {Ball velocity and beam-angle rate estimates derived from the logged telemetry.}
  {fig:telemetry_velocity_rate}

The logged beam-angle command remains well inside the configured \SI{2.0}{deg} command
cap, and the CSV shows \verb|theta_cmd_sat = 0| for all rows. Recovery activates only
briefly during large transitions, which is consistent with the firmware design intent:
use recovery as a minimum-corrective-angle floor rather than as the dominant transport
mechanism.

\section{Performance Metrics}

\ReportTableFromFile
  {Per-step tracking metrics computed from valid closed-loop samples. Settling uses the
  default band $\max(\SI{0.2}{cm}, 5\%$ of commanded step size$)$. Unsupported metrics
  remain \texttt{N/A}.}
  {tables/telemetry_step_metrics.tex}
  {tab:telemetry_step_metrics}

The step metrics confirm three main points:

\begin{itemize}
  \item the far-side transition is limited by invalid distance samples and does not
  reach a clean settled condition in the available segment length,
  \item the center-return and near-side transitions achieve finite rise times but still
  show non-negligible residual error over the final valid window, and
  \item the recorded experiment is more useful as a realistic tuning baseline than as a
  final high-performance demonstration.
\end{itemize}
